% GENERATED FILE. Edit canonical lessons, then run npm run generate:books.
% canonical-source-hash: fnv1a64:0236bf6ca1b0dbd0
% canonical-lessons: BN-C432-chura, BN-C432-upotyoka, BN-C432-dhibi, BN-C432-morubhumi, BN-C432-jonggol

\chapter{Peaks, Valleys and Deserts}
\label{ch:bn-peaks-valleys-and-deserts}
\hlchaptermodality{432}

\begin{chapteropening}
\textbf{By the end of this chapter:} \emph{I can describe a summit, a valley, a mound, a desert and a forest.}

Complete the last lesson of chapter 432: I can describe a summit, a valley, a mound, a desert and a forest.
\end{chapteropening}

\section[\texorpdfstring{\bn{চূড়া} (chūṛā)}{চূড়া (chūṛā)}]{\bn{চূড়া} (chūṛā) — a peak; a summit; a top}

\label{lesson:BN-C432-chura}



\begin{quote}
Before the new one: say the Bengali for silt, then the Bengali for to fill in (a pond or low land).
\end{quote}

\subsection*{You'll want to know: \bn{চূড়া}}
\textbf{\bn{চূড়া}} — \emph{chūṛā} — \textquotedblleft{}a peak; a summit; a top\textquotedblright{}.

\begin{quote}
\bn{আমরা} \bn{ভোরে} \bn{পাহাড়ের} \bn{চূড়ায়} \bn{উঠলাম।}
\end{quote}

\begin{quote}
\emph{āmrā bhore pāhāṛer chūṛāy uṭhlām.}
\end{quote}

\begin{quote}
We climbed to the summit at dawn.
\end{quote}

Also the top of anything tall: \bn{মন্দিরের} \bn{চূড়া}, the spire of a temple.

\begin{cousinweb}
\textbf{\bn{চূড়া}} (\emph{chūṛā}) is Sanskrit \emph{cūḍā}, \textquotedblleft{}the hair on the top of the head, a crest\textquotedblright{} (Monier-Williams); the sense of a mountain top is a natural step, though Monier-Williams gives the head first.
\end{cousinweb}

\begin{grammarlens}[title={What you've built}]
One more word for getting around.
\end{grammarlens}

\subsection*{Your turn}
\begin{itemize}
\raggedright
  \item \emph{Say it:} \emph{chūṛā}
  \item \emph{Say it:} \emph{chūṛā}, once more
  \item \emph{Say it:} \emph{āmrā bhore pāhāṛer chūṛāy uṭhlām.} — We climbed to the summit at dawn.
  \item \emph{Recall:} say the Bengali for silt, then the Bengali for to fill in (a pond or low land), then say \emph{chūṛā} again
\end{itemize}

\begin{tcolorbox}[breakable,colback=teal!4,colframe=teal!35!black,title={Before you move on}]
What does \bn{চূড়া} mean? (\textquotedblleft{}A peak; a summit; a top\textquotedblright{}.) Say it once more.
\end{tcolorbox}
\section[\texorpdfstring{\bn{উপত্যকা} (upôtyôkā)}{উপত্যকা (upôtyôkā)}]{\bn{উপত্যকা} (upôtyôkā) — a valley}

\label{lesson:BN-C432-upotyoka}



\begin{quote}
Before the new one: say the Bengali for to fill in (a pond or low land), then the Bengali for a summit.
\end{quote}

\subsection*{You'll want to know: \bn{উপত্যকা}}
\textbf{\bn{উপত্যকা}} — \emph{upôtyôkā} — \textquotedblleft{}a valley\textquotedblright{}.

\begin{quote}
\bn{পাহাড়ের} \bn{নিচে} \bn{উপত্যকায়} \bn{ছোট} \bn{একটা} \bn{গ্রাম।}
\end{quote}

\begin{quote}
\emph{pāhāṛer niche upôtyôkāy chhoṭo ekṭā grām.}
\end{quote}

\begin{quote}
Below the mountain, in the valley, lies a small village.
\end{quote}

\begin{cousinweb}
\textbf{\bn{উপত্যকা}} (\emph{upôtyôkā}) is Sanskrit \emph{upatyakā}, \textquotedblleft{}land at the foot of a mountain, low land\textquotedblright{} (Monier-Williams), taken over whole.
\end{cousinweb}

\begin{grammarlens}[title={What you've built}]
One more word for getting around.
\end{grammarlens}

\subsection*{Your turn}
\begin{itemize}
\raggedright
  \item \emph{Say it:} \emph{upôtyôkā}
  \item \emph{Say it:} \emph{upôtyôkā}, once more
  \item \emph{Say it:} \emph{pāhāṛer niche upôtyôkāy chhoṭo ekṭā grām.} — Below the mountain, in the valley, lies a small village.
  \item \emph{Recall:} say the Bengali for to fill in (a pond or low land), then the Bengali for a summit, then say \emph{upôtyôkā} again
\end{itemize}

\begin{tcolorbox}[breakable,colback=teal!4,colframe=teal!35!black,title={Before you move on}]
What does \bn{উপত্যকা} mean? (\textquotedblleft{}A valley\textquotedblright{}.) Say it once more.
\end{tcolorbox}
\section[\texorpdfstring{\bn{ঢিবি} (ḍhibi)}{ঢিবি (ḍhibi)}]{\bn{ঢিবি} (ḍhibi) — a mound; a heap of earth}

\label{lesson:BN-C432-dhibi}



\begin{quote}
Before the new one: say the Bengali for a summit, then the Bengali for a valley.
\end{quote}

\subsection*{You'll want to know: \bn{ঢিবি}}
\textbf{\bn{ঢিবি}} — \emph{ḍhibi} — \textquotedblleft{}a mound; a heap of earth\textquotedblright{}.

\begin{quote}
\bn{মাঠের} \bn{মাঝখানে} \bn{একটা} \bn{মাটির} \bn{ঢিবি} \bn{আছে।}
\end{quote}

\begin{quote}
\emph{māṭher mājhkhāne ekṭā māṭir ḍhibi āchhe.}
\end{quote}

\begin{quote}
There is a mound of earth in the middle of the field.
\end{quote}

\begin{cousinweb}
\textbf{\bn{ঢিবি}} (\emph{ḍhibi}) is inherited: CDIAL lists Bengali \emph{ḍhibi}, \textquotedblleft{}a heap\textquotedblright{}, under the reconstructed \emph{ḍhappa}, \textquotedblleft{}a lump\textquotedblright{} (CDIAL 5580).
\end{cousinweb}

\begin{grammarlens}[title={What you've built}]
One more word for getting around.
\end{grammarlens}

\subsection*{Your turn}
\begin{itemize}
\raggedright
  \item \emph{Say it:} \emph{ḍhibi}
  \item \emph{Say it:} \emph{ḍhibi}, once more
  \item \emph{Say it:} \emph{māṭher mājhkhāne ekṭā māṭir ḍhibi āchhe.} — There is a mound of earth in the middle of the field.
  \item \emph{Recall:} say the Bengali for a summit, then the Bengali for a valley, then say \emph{ḍhibi} again
\end{itemize}

\begin{tcolorbox}[breakable,colback=teal!4,colframe=teal!35!black,title={Before you move on}]
What does \bn{ঢিবি} mean? (\textquotedblleft{}A mound; a heap of earth\textquotedblright{}.) Say it once more.
\end{tcolorbox}
\section[\texorpdfstring{\bn{মরুভূমি} (môrubhūmi)}{মরুভূমি (môrubhūmi)}]{\bn{মরুভূমি} (môrubhūmi) — a desert}

\label{lesson:BN-C432-morubhumi}



\begin{quote}
Before the new one: say the Bengali for a valley, then the Bengali for a mound.
\end{quote}

\subsection*{You'll want to know: \bn{মরুভূমি}}
\textbf{\bn{মরুভূমি}} — \emph{môrubhūmi} — \textquotedblleft{}a desert\textquotedblright{}.

\begin{quote}
\bn{মরুভূমিতে} \bn{দিনে} \bn{গরম}, \bn{রাতে} \bn{ঠান্ডা।}
\end{quote}

\begin{quote}
\emph{môrubhūmite dine gôrôm, rāte ṭhānḍā.}
\end{quote}

\begin{quote}
In the desert it is hot by day and cold by night.
\end{quote}

\begin{cousinweb}
\textbf{\bn{মরুভূমি}} (\emph{môrubhūmi}) joins Sanskrit \emph{maru}, \textquotedblleft{}a wilderness, a sandy waste\textquotedblright{}, to \emph{bhūmi}, \textquotedblleft{}land, ground\textquotedblright{} (Monier-Williams for both).
\end{cousinweb}

\begin{grammarlens}[title={What you've built}]
One more word for getting around.
\end{grammarlens}

\subsection*{Your turn}
\begin{itemize}
\raggedright
  \item \emph{Say it:} \emph{môrubhūmi}
  \item \emph{Say it:} \emph{môrubhūmi}, once more
  \item \emph{Say it:} \emph{môrubhūmite dine gôrôm, rāte ṭhānḍā.} — In the desert it is hot by day and cold by night.
  \item \emph{Recall:} say the Bengali for a valley, then the Bengali for a mound, then say \emph{môrubhūmi} again
\end{itemize}

\begin{tcolorbox}[breakable,colback=teal!4,colframe=teal!35!black,title={Before you move on}]
What does \bn{মরুভূমি} mean? (\textquotedblleft{}A desert\textquotedblright{}.) Say it once more.
\end{tcolorbox}
\section[\texorpdfstring{\bn{জঙ্গল} (jônggôl)}{জঙ্গল (jônggôl)}]{\bn{জঙ্গল} (jônggôl) — a forest; a jungle; overgrown land}

\label{lesson:BN-C432-jonggol}



\begin{quote}
Before the new one: say the Bengali for a mound, then the Bengali for a desert.
\end{quote}

\subsection*{You'll want to know: \bn{জঙ্গল}}
\textbf{\bn{জঙ্গল}} — \emph{jônggôl} — \textquotedblleft{}a forest; a jungle; overgrown land\textquotedblright{}.

\begin{quote}
\bn{গ্রামের} \bn{পিছনে} \bn{একটা} \bn{বড়} \bn{জঙ্গল} \bn{আছে।}
\end{quote}

\begin{quote}
\emph{grāmer pichhône ekṭā bôṛo jônggôl āchhe.}
\end{quote}

\begin{quote}
Behind the village there is a big forest.
\end{quote}

\bn{বন} is a forest in general; \bn{জঙ্গল} is wild and thick, and also an overgrown patch of weeds.

\begin{cousinweb}
\textbf{\bn{জঙ্গল}} (\emph{jônggôl}) is inherited: CDIAL lists Bengali \emph{jaṅgal}, \textquotedblleft{}a forest\textquotedblright{}, under Sanskrit \emph{jāṅgala}, \textquotedblleft{}arid, sparsely wooded; a desert region\textquotedblright{} (CDIAL 5177); the turn from dry land to thick forest happened in the modern languages, and English \emph{jungle} is the same word.
\end{cousinweb}

\begin{grammarlens}[title={What you've built}]
One more word for getting around.
\end{grammarlens}

\subsection*{Your turn}
\begin{itemize}
\raggedright
  \item \emph{Say it:} \emph{jônggôl}
  \item \emph{Say it:} \emph{jônggôl}, once more
  \item \emph{Say it:} \emph{grāmer pichhône ekṭā bôṛo jônggôl āchhe.} — Behind the village there is a big forest.
  \item \emph{Recall:} say the Bengali for a mound, then the Bengali for a desert, then say \emph{jônggôl} again
\end{itemize}

\begin{tcolorbox}[breakable,colback=teal!4,colframe=teal!35!black,title={Before you move on}]
What does \bn{জঙ্গল} mean? (\textquotedblleft{}A forest; a jungle; overgrown land\textquotedblright{}.) Say it once more.
\end{tcolorbox}
